\documentclass[a4paper]{article}
% Español (México) con XeLaTeX: fontspec para fuentes Unicode, polyglossia
% para la división silábica y los nombres del documento en la variante mexicana.
\usepackage{fontspec}
\usepackage{polyglossia}
\setdefaultlanguage[variant=mexican]{spanish}
% Los nombres chinos van como «pinyin (original)»: xeCJK da a los caracteres
% Han su propia fuente; sin ella XeLaTeX los omite con un aviso.
% FandolSong no cubre algunos caracteres de nombres (U+73FA); WenQuanYi Zen Hei sí.
\usepackage[AutoFallBack=true]{xeCJK}
\setCJKmainfont{FandolSong-Regular.otf}
\setCJKfallbackfamilyfont{\CJKrmdefault}{WenQuanYi Zen Hei}
% polyglossia no traduce los nombres de listings.
\AtBeginDocument{\providecommand{\lstlistingname}{}\renewcommand{\lstlistingname}{Listado}%
  \providecommand{\lstlistlistingname}{}\renewcommand{\lstlistlistingname}{Índice de listados}}
\usepackage{amsmath, amssymb}
\usepackage{graphicx}
\usepackage[margin=2.5cm]{geometry}
\usepackage[most]{tcolorbox}
\usepackage{etoolbox}
\usepackage{listings}
\usepackage{hyperref}
\usepackage{booktabs}
\usepackage{subcaption}
\usepackage{float}

\newtcolorbox{knowledgebox}[1]{enhanced, colback=blue!5!white, colframe=blue!75!black, colbacktitle=blue!75!black, coltitle=white, fonttitle=\bfseries, title={#1}, attach boxed title to top left={yshift=-2mm, xshift=2mm}, boxrule=1pt, sharp corners}
\newtcolorbox{importantbox}[1]{enhanced, colback=yellow!10!white, colframe=yellow!80!black, colbacktitle=yellow!80!black, coltitle=black, fonttitle=\bfseries, title={#1}, sharp corners}
\newtcolorbox{warningbox}[1]{enhanced, colback=red!5!white, colframe=red!75!black, colbacktitle=red!75!black, coltitle=white, fonttitle=\bfseries, title={#1}, sharp corners}

\lstset{language=Python, basicstyle=\ttfamily\small, keywordstyle=\color{blue}, stringstyle=\color{red!60!black}, commentstyle=\color{green!60!black}, breaklines=true, frame=single, numbers=left, numberstyle=\tiny\color{gray}, captionpos=b, extendedchars=true}

\newcommand{\notetitle}{{[}LLM Agents F24{]} Agentic AI Frameworks \& AutoGen + Building a Multimodal Knowledge Assistant}
\newcommand{\noteauthors}{Elaboradas a partir de materiales públicos del curso}
\newcommand{\notedate}{23 de septiembre de 2024}
\newcommand{\videochannel}{Berkeley RDI}
\newcommand{\videopublishdate}{23 de septiembre de 2024}
\newcommand{\videoduration}{aproximadamente 80 minutos}
\newcommand{\videourl}{https://www.youtube.com/watch?v=OOdtmCMSOo4}
\newcommand{\videocoverpath}{cover.jpg}
\newcommand{\repourl}{https://github.com/hqhq1025/ai-course-notes}


% Footer with repo info
\usepackage{fancyhdr}
\pagestyle{fancy}
\fancyhf{}
\fancyfoot[C]{\thepage}
\fancyfoot[R]{\footnotesize\color{gray}\href{\repourl}{AI Course Notes}}
\renewcommand{\headrulewidth}{0pt}
\renewcommand{\footrulewidth}{0pt}

\begin{document}
\begin{titlepage}\centering
{\Large Notas del curso\par}\vspace{1.2cm}
{\huge\bfseries \notetitle\par}\vspace{0.8cm}
{\large \noteauthors\par}\vspace{0.3cm}
{\large \notedate\par}\vspace{1.2cm}
\ifdefempty{\videocoverpath}{{\small Al generar las notas, indique la ruta local de la portada del video.\par}}{\includegraphics[width=0.82\textwidth,height=0.45\textheight,keepaspectratio]{\videocoverpath}\par}
\vfill
\begin{tcolorbox}[width=0.9\textwidth, colback=black!2!white, colframe=black!60, sharp corners]
\textbf{Autor o canal del video}: \videochannel\par
\textbf{Fecha de publicación}: \videopublishdate\par
\textbf{Duración del video}: \videoduration\par
\ifdefempty{\videourl}{}{\textbf{Enlace del video}: \href{\videourl}{\nolinkurl{\videourl}}\par}\par
\textbf{Página del proyecto}: \href{\repourl}{\nolinkurl{\repourl}}\par
\end{tcolorbox}
\end{titlepage}

\tableofcontents\newpage

\section{Introducción: el futuro de la IA agéntica}

Esta clase la presentan en conjunto dos invitados: \textbf{Chi Wang} (Microsoft Research, desarrollador principal del framework AutoGen) presenta los frameworks de IA agéntica y AutoGen; \textbf{Jerry Liu} (fundador y director general de LlamaIndex) presenta la construcción de asistentes de conocimiento multimodales.

\subsection{¿Qué es la IA agéntica?}

\begin{importantbox}{El valor central de la IA agéntica}
\begin{enumerate}
    \item \textbf{Interfaz en lenguaje natural}: el usuario describe sus necesidades en lenguaje natural y el agente las entiende y ejecuta automáticamente
    \item \textbf{Ejecución autónoma de tareas complejas}: mediante razonamiento de varios pasos y llamadas a herramientas, completa tareas con la mínima supervisión humana
    \item \textbf{Nueva arquitectura de software}: varios agentes colaboran para formar un sistema modular y componible de manera recursiva
\end{enumerate}
\end{importantbox}

El ejemplo muestra el proceso completo en que Magentic-One (basado en AutoGen) construye un sitio web de forma automática: analiza la tarea, instala las dependencias, varios agentes colaboran para escribir el código, lo compila y lo ejecuta, e incluso es capaz de autorrepararse después de que se elimina código clave.

\subsection{Resumen de la sección}

La IA agéntica no es solo una capa de aplicación de los LLM, sino que representa un paradigma completamente nuevo de desarrollo de software: pasar de «escribir código» a «orquestar agentes».

\section{AutoGen: framework de programación de conversaciones entre múltiples Agentes}

\subsection{Filosofía central de diseño}

El diseño de AutoGen se basa en una intuición central: \textbf{la conversación (conversation) es el primitivo más básico de interacción entre Agentes}.

\begin{knowledgebox}{Por qué elegir la conversación como primitivo básico}
La conversación es la forma básica en que los seres humanos avanzan en el aprendizaje y demuestran teoremas. El éxito de ChatGPT muestra que las conversaciones de varios turnos se adaptan de manera natural a la mejora iterativa y a la resolución colaborativa. Un diseño centrado en la conversación puede derivar de forma natural todos los demás patrones de interacción entre Agentes.
\end{knowledgebox}

Construir una aplicación de Agentes solo requiere dos pasos:
\begin{enumerate}
    \item \textbf{Definir los Agentes}: cada Agente puede usar un LLM, herramientas o entrada humana como backend
    \item \textbf{Hacer que los Agentes conversen}: orquestar la colaboración mediante patrones de conversación (conversación entre dos, conversación anidada, chat grupal, etc.)
\end{enumerate}

\subsection{La abstracción Conversable Agent}

La abstracción unificada de AutoGen, \textbf{Conversable Agent}, unifica distintos tipos de entidades (LLM, herramientas, seres humanos) como Agentes capaces de conversar:
\begin{itemize}
    \item Se puede elegir un LLM como backend (por ejemplo, GPT-4, Claude)
    \item Se pueden usar herramientas (ejecución de código Python, llamadas a API, etc.)
    \item Se puede conectar entrada humana
    \item Se pueden combinar varios backends
    \item Admite anidamiento: un Agente puede contener en su interior conversaciones de otros Agentes
\end{itemize}

\subsection{Patrones de conversación}

\begin{itemize}
    \item \textbf{Conversación entre dos}: el patrón más simple, ya admite Reflection (un Agente escritor + un Agente revisor que mejoran de forma iterativa)
    \item \textbf{Sequential Chat}: varias conversaciones entre dos encadenadas, como revisiones anidadas desde múltiples perspectivas (revisión de SEO → revisión legal → revisión ética → metarrevisión)
    \item \textbf{Nested Chat}: un Agente inicia subconversaciones en su interior, como cuando el Commander Agent primero conversa con el Writer para obtener el código y después conversa con el Safeguard para verificar su seguridad
    \item \textbf{Group Chat}: varios Agentes con roles distintos eligen dinámicamente al que habla dentro de un grupo, con soporte para condiciones de restricción y lógica de transición de estados
\end{itemize}

\subsection{Ventajas de Multi-Agent}

\begin{importantbox}{¿Por qué múltiples Agentes son mejores que uno solo?}
Los experimentos muestran que descomponer una tarea entre Agentes especializados (como Writer + Safeguard) produce mejores resultados que hacer que un solo Agente atienda todas las instrucciones a la vez: el esquema de múltiples Agentes con GPT-4 supera al esquema de un solo Agente en 20\% en las verificaciones de seguridad, y con GPT-3.5 la diferencia es aún mayor. Cuanto más compleja es la tarea y más débil el modelo, más evidente es la ventaja de múltiples Agentes.
\end{importantbox}

Desde el punto de vista de la programación, la arquitectura de múltiples Agentes también ofrece:
\begin{itemize}
    \item Mejor modularidad y mantenibilidad
    \item Soporte natural para la participación humana (un humano puede tomar el lugar de cualquier rol de Agente)
    \item Facilidad de extensión (reemplazar un solo Agente no afecta al conjunto)
\end{itemize}

\subsection{AutoBuild: construcción automática de equipos de Agentes}

AutoBuild resuelve el problema de ``qué equipo de Agentes diseñar para una tarea específica'':
\begin{enumerate}
    \item El usuario describe el requerimiento de alto nivel
    \item El sistema sugiere automáticamente Agentes con distintos roles
    \item Los Agentes forman un Group Chat para colaborar en la tarea
    \item Adaptive Build además admite construir el equipo de Agentes de forma dinámica y por pasos
\end{enumerate}

\subsection{Resumen de la sección}

AutoGen toma la conversación como primitivo central y ofrece un framework de programación de múltiples Agentes flexible y potente. Su filosofía de diseño es ``definir Agentes + hacer que conversen''; mediante el anidamiento y la composición, es posible construir de manera recursiva sistemas de Agentes de complejidad arbitraria.

\section{LlamaIndex: construcción de un asistente de conocimiento multimodal}

\subsection{De Basic RAG a Advanced RAG}

Jerry Liu señala cuatro limitaciones fundamentales del Basic RAG:
\begin{enumerate}
    \item \textbf{Procesamiento de datos tosco}: la simple división del texto en chunks no puede manejar tablas, imágenes ni diseños complejos
    \item \textbf{El LLM se usa solo para la síntesis}: no se aprovechan las capacidades de razonamiento y planeación
    \item \textbf{Sin estado}: cada consulta es independiente, sin memoria personalizada
    \item \textbf{Salida única}: solo puede generar respuestas de texto simples
\end{enumerate}

\begin{warningbox}{Garbage In Garbage Out}
La calidad de una aplicación de RAG depende de la calidad de su pipeline de procesamiento de datos. Si el analizador de PDF no puede extraer correctamente las tablas y las imágenes, ningún razonamiento posterior del LLM, por potente que sea, puede compensarlo. La calidad de los datos es una condición necesaria para las aplicaciones de LLM en producción.
\end{warningbox}

\subsection{Pipeline de RAG multimodal}

Pasos clave para construir un RAG multimodal:

\begin{enumerate}
    \item \textbf{Análisis de documentos de alta calidad}: extraer por separado el texto, las tablas, las gráficas y las imágenes de un PDF como elementos estructurados
    \item \textbf{Indexación jerárquica}: generar para cada elemento múltiples representaciones textuales (resúmenes, sub-chunks, etc.) que apunten al elemento fuente
    \item \textbf{Recuperación multimodal}: dada una consulta, recuperar chunks de texto e imagen, aprovechando un LLM multimodal para procesar simultáneamente la información textual y la visual
\end{enumerate}

\begin{knowledgebox}{Nuevas oportunidades de los modelos multimodales}
Los modelos dominantes actuales (GPT-4o, Claude 3.5 Sonnet, Gemini) admiten la entrada simultánea de texto e imágenes. Esto significa que, incluso si la indexación se construye con un modelo de embeddings de texto, en la etapa final de inferencia aún es posible alimentar la imagen original al LLM y lograr una comprensión visual genuina.
\end{knowledgebox}

\subsection{Agentic RAG: agregar una capa de razonamiento sobre el RAG}

\begin{importantbox}{La idea central del Agentic RAG}
Tratar la recuperación de RAG como una herramienta (tool) y agregar por encima una capa de razonamiento de Agent. El Agent puede decidir cómo descomponer el problema, qué herramientas de recuperación invocar y cuándo hacer una reflexión de verificación, de modo que pueda abordar problemas complejos que el RAG básico no puede responder (resúmenes, comparación de múltiples documentos, análisis de varios pasos, etc.).
\end{importantbox}

Dos opciones de arquitectura:
\begin{itemize}
    \item \textbf{Flujo restringido}: el desarrollador define manualmente la lógica if-else y las reglas de enrutamiento; es más confiable pero menos flexible
    \item \textbf{Arquitectura de Agent general} (como ReAct): el Agent elige por sí mismo las herramientas y la ruta de razonamiento; es más flexible pero menos confiable y de mayor costo
\end{itemize}

Regla práctica: en una arquitectura de Agent general, el número de herramientas debe mantenerse en lo posible entre 4 y 5, sin superar las 10.

\subsection{Despliegue en producción}

LlamaIndex ofrece un sistema de orquestación basado en eventos que permite:
\begin{itemize}
    \item Encapsular cada flujo de trabajo de Agent como una API de microservicio
    \item Comunicación entre Agents mediante una cola de mensajes central
    \item Human-in-the-loop: el Agent puede pausarse para esperar la entrada de un humano antes de continuar la ejecución
    \item Soporte para desarrollo local y despliegue en Kubernetes
\end{itemize}

\subsection{Resumen de la sección}

LlamaIndex, partiendo de la calidad del procesamiento de datos, construye un pipeline completo que va del análisis de documentos a la recuperación multimodal y de ahí al razonamiento Agentic. Su idea central es: una buena calidad de datos es la base, el razonamiento del Agent es la superestructura, y ninguna de las dos puede faltar.

\section{Síntesis y ampliación}

\subsection{Puntos clave}

\begin{enumerate}
    \item \textbf{AutoGen} toma la conversación como abstracción central, admite la composición recursiva de múltiples Agentes y es adecuado para construir sistemas complejos de colaboración entre Agentes.
    \item \textbf{LlamaIndex} parte de la calidad de los datos y construye asistentes de conocimiento de nivel de producción mediante RAG multimodal y razonamiento Agentic.
    \item Las arquitecturas de múltiples Agentes tienen ventajas notables frente a un solo Agente en la descomposición de tareas, la garantía de seguridad y la capacidad de mantenimiento.
    \item Entre las arquitecturas de Agente restringidas y las de propósito general existe un compromiso entre confiabilidad y flexibilidad.
    \item El despliegue en producción debe considerar la conversión a microservicios, las colas de mensajes y human-in-the-loop.
\end{enumerate}

\subsection{Lecturas adicionales}

\begin{itemize}
    \item Wu et al., ``AutoGen: Enabling Next-Gen LLM Applications via Multi-Agent Conversation,'' 2023.
    \item Liu, ``LlamaIndex: A Data Framework for LLM Applications,'' 2022--2024.
    \item Lewis et al., ``Retrieval-Augmented Generation for Knowledge-Intensive NLP Tasks,'' NeurIPS 2020.
\end{itemize}

\end{document}
